\section{Need-Sufficiency Architecture}
\label{sec:need_sufficiency}

The revealed-sacrifice observation channel
(Section~\ref{sec:sacrifice_model}) assumes all sacrifices have
the same polarity: more sacrifice equals more revealed value.
This is wrong for a class of sacrifices that are structurally
involuntary.  An agent who pays \$500 for insulin is not revealing
that insulin is worth \$500 beyond baseline, they are revealing
that baseline functionality costs \$500 to maintain.  The sacrifice
magnitude reveals \emph{cost-of-access} (how expensive it is for
this agent to reach minimum viable functionality), not
\emph{preference strength} (how much the agent values the outcome
beyond baseline).  Larger sacrifice means a worse situation, not
higher revealed value.

The framework needs a structural requirement on poset
construction that distinguishes need-satisfaction from want-pursuit
and identifies the boundary where sacrifice polarity flips.

\subsection{Three motivating observations}

\paragraph{Needs are neither purely maximizing nor minimizing.}
Every need has \emph{both} maximizing dimensions (quality, quantity,
reliability) and minimizing dimensions (cost, distance, time). Needs are not a
single-polarity class. What \emph{is} true: the minimizing dimensions of needs
are where the polarity problem lives, because cost-of-access is involuntary.

\paragraph{Scalar sufficiency hides structural defects.}
A single ``water access: sufficient'' measurement might hide that the water is
contaminated.  The agent has water (passes the scalar threshold) but the way
the need is met actively degrades future $\volL$ through downstream effects on
health, cognition, and reproduction.  The framework reports green while the
system corrodes from the root.

\paragraph{Downstream-safe bundles.}
Three levels of sufficiency, each strictly stronger:
\begin{itemize}
\item \emph{Level~1: Scalar threshold (insufficient).}
  ``Water: yes/no.''  Hides quality, reliability, and cost defects.
\item \emph{Level~2: Bundle threshold (necessary but not
  sufficient).}
  All dimensions of the need must independently meet their thresholds
        simultaneously. This catches contamination (quality dimension fails),
        but misses subtler failures where the satisfaction pathway degrades
        downstream capabilities over time. For example, when the contamination
        thresholds are incorrectly tuned.
\item \emph{Level~3: Downstream-safe bundle.}
  Bundle thresholds are all met \emph{and} the need-satisfaction pathway does
        not produce net-negative effects on capabilities in the need's
        downstream capability cone.
\end{itemize}

\subsection{Formal definitions}

\begin{definition}[Multi-Dimensional Need Bundle]
\label{def:need_bundle}
A capability $N$ identified as a need at position $p_N$ in the
poset has:
\begin{itemize}
\item \textbf{Dimensions} $\{d_1, \ldots, d_m\}$ with sufficiency
  thresholds $\{s_1, \ldots, s_m\}$.
\item Each dimension has a \textbf{polarity}: maximizing (quality,
  quantity, reliability) or minimizing (cost, distance, time).
\item \textbf{Downstream cone}
  $\Down(p_N) = \{c : c \text{ depends on } N \text{ in the
  poset}\}$.
\item \textbf{Satisfaction pathway} $\pi_N$: the specific way the
  agent meets the need (which source, at what quality, at what
  cost).
\end{itemize}
For cost dimensions (minimizing polarity), the benchmark inverts
the natural direction: $\mathrm{actual}_k = 1 / \mathrm{cost}_k$,
with sufficiency at $s_k = 1 / \mathrm{max\_acceptable\_cost}$,
so that benchmark maximization aligns with cost minimization.
\end{definition}

\begin{definition}[Downstream-Safe Sufficiency]
\label{def:downstream_safe}
Need $N$ at position $p_N$ is \emph{downstream-safe sufficient}
under satisfaction pathway $\pi_N$ iff:
\begin{enumerate}
\item[(a)] \textbf{Bundle thresholds met:}
  $d_k(\pi_N) \geq s_k$ for all dimensions~$k$ (with cost
  dimensions inverted per Definition~\ref{def:need_bundle}).
\item[(b)] \textbf{Downstream non-degradation:} For every
  capability $c \in \Down(p_N)$, the quality of $c$ under the
  actual satisfaction pathway $\pi_N$ is at least as high as the
  threshold-reference quality $q_N^0(c)$:
  \[
    \mathrm{quality}_c(\text{with } N \text{ satisfied via }
    \pi_N) \;\geq\; q_N^0(c),
  \]
  where $q_N^0(c)$ is the pointwise threshold-reference quality
  defined below.
\end{enumerate}
\end{definition}

\begin{definition}[Threshold-Reference Quality]
\label{def:threshold_ref}
Let $\Pi_N^0 = \{\pi : d_k(\pi) = s_k \text{ for all } k\}$ be
the class of pathways satisfying each dimension of~$N$ at exactly
its sufficiency threshold.

\emph{Preconditions:} (i)~$\Pi_N^0$ is nonempty---at least one
pathway exists that satisfies every dimension at exactly the
threshold level.  (ii)~For each $c \in \Down(p_N)$, the function
$\pi \mapsto \mathrm{quality}_c(\text{with } N \text{ satisfied
via } \pi, \text{ all other inputs held fixed})$ is bounded above
on $\Pi_N^0$, so the supremum exists in $\mathbb{R}$.

For each downstream capability $c \in \Down(p_N)$:
\[
  q_N^0(c) \;=\; \sup_{\pi \in \Pi_N^0}
  \mathrm{quality}_c(\text{with } N \text{ satisfied via } \pi).
\]
The pointwise supremum avoids requiring a single pathway that
simultaneously maximizes downstream quality for all
$c \in \Down(p_N)$.  $q_N^0(c)$ is a \emph{definitional
construct}, not an observed quantity---it serves as the most
favourable threshold-level baseline for each downstream capability
separately.
\end{definition}

\begin{remark}[Structural connection to exercise-pathway machinery]
\label{rem:exercise_pathway}
The need-satisfaction pathway $\pi_N$ is \emph{structurally} an
exercise pathway in the sense
of~\citep[Definition~2]{lasser2026horizon}: a sequence of
capability invocations producing an outcome (here, the
satisfaction of need~$N$).  This identification lets the framework
treat $\pi_N$ as a first-class object rather than an abstract
``way of satisfying the need.''

The downstream-safety condition
(Definition~\ref{def:downstream_safe}(b)) is, however, a distinct
mathematical object from the concentration-risk measure
of~\citep[Definition~3]{lasser2026horizon}.  Concentration risk
is a scalar worst-case contraction estimate over an entire
pathway; downstream-safety is a per-capability quality
comparison $\mathrm{quality}_c(\pi_N) \geq q_N^0(c)$ indexed by
$c \in \Down(p_N)$.  Neither reduces to the other.
Evaluating $\mathrm{quality}_c(\pi_N)$ in practice requires
either (i) direct observation of $c$'s performance conditional on
the pathway, or (ii) structural inference from the poset's
cooperative structure.  A full formalisation of the evaluation
procedure is outside this paper's scope; within the paper,
Definition~\ref{def:downstream_safe}(b) is a structural condition
the framework \emph{enforces} (a pathway failing the condition is
rejected as a satisfaction pathway for $N$), not a quantity the
framework \emph{computes} directly.

A policy consequence, not a theorem: because need-satisfaction
pathways are not optional (the agent must satisfy the need), the
framework's response to a pathway failing
Definition~\ref{def:downstream_safe}(b) differs from its response
to a risk claim under~\citep{lasser2026wamura}.  A failing risk
pathway is \emph{restricted} (the capability is blocked).  A
failing satisfaction pathway must be \emph{replaced} (the need
must still be met, but through a pathway satisfying
Definition~\ref{def:downstream_safe}).  This governance
distinction is upstream of the paper's formal machinery; it names
the correct operational response when the condition fails, not a
new computation the framework performs.
\end{remark}

\begin{definition}[Sufficiency Threshold as Polarity Boundary]
\label{def:polarity_boundary}
The sufficiency bundle is the structural boundary where
revealed-sacrifice polarity flips:
\begin{itemize}
\item \textbf{Below sufficiency on any dimension:} The agent is
  investing to reach minimum viable functionality.  Sacrifice is
  involuntary (the agent must reach sufficiency to operate).
  Sacrifice magnitude reveals \textbf{cost-of-access}---how
  expensive it is for this agent to reach baseline on this
  dimension.  Larger sacrifice = worse situation.
\item \textbf{Above sufficiency on all dimensions:} The agent has
  reached functionality and is choosing to invest further.
  Sacrifice is voluntary (the agent could stop at ``good
  enough'').  Sacrifice magnitude reveals \textbf{preference
  strength}---how much the agent values improvement beyond
  baseline.  Larger sacrifice = higher revealed value.
\end{itemize}
The sufficiency threshold is where the need-insufficiency failure
mode of Assumption~S1 (Section~\ref{sec:sacrifice_model}) is
removed: below sufficiency, S1 fails because retaining the
sacrificed capability is not a non-degrading alternative; above
sufficiency, this specific failure mode does not apply (though S1
can still fail via other mechanisms---see
Proposition~\ref{prop:s1_sufficiency}, Remark).
\end{definition}

\begin{definition}[Polarity-Correct Benchmark]
\label{def:polarity_benchmark}
For each dimension $d_k$ of a need capability:
\begin{itemize}
\item \textbf{Below sufficiency:}
  $\mathrm{benchmark}_k = \min(\mathrm{actual}_k / s_k,\, 1)$.
  Climbs from $0$ toward $1$ as the agent approaches sufficiency.
  Caps at $1$.  Maximizing this equals closing the sufficiency
  gap.
\item \textbf{Above sufficiency:}
  $\mathrm{benchmark}_k = 1 + \alpha_k \cdot
  \max\!\bigl((\mathrm{actual}_k - s_k)/s_k,\, 0\bigr)$,
  where $\alpha_k < 1$ discounts above-sufficiency gains relative
  to below-sufficiency gains.  This makes closing the sufficiency
  gap more valuable per unit of improvement than enhancing beyond
  sufficiency---the correct incentive structure for needs.
\end{itemize}
\end{definition}

\subsection{Structural propositions}

\begin{proposition}[Anti-Stuffing]
\label{prop:anti_stuffing}
Under downstream-safe bundle sufficiency
(Definition~\ref{def:downstream_safe}), redundant capability
measurements---e.g., ``distance from $A$ to gas station,''
``distance from $B$ to gas station,'' measuring the same need
from different vantage points---provide no downstream-quality
improvement.  Once the bundle is sufficient, the
below-sufficiency benchmark caps at $1$
(Definition~\ref{def:polarity_benchmark}), and the
above-sufficiency component credits only genuine quality
improvement (cheaper fuel, faster access), not re-measurement
from different positions.  The bundle structure plus
downstream-safety naturally resists stuffing without a separate
anti-stuffing mechanism.
\end{proposition}

\begin{proof}[Proof sketch]
Under Definition~\ref{def:downstream_safe}, sufficiency is a
property of the \emph{bundle} (all dimensions simultaneously),
not of individual measurements.  A stuffed measurement
contributes to a dimension already at or above threshold.  By
the capping discipline (Definition~\ref{def:polarity_benchmark}),
its below-sufficiency benchmark is
$\min(\mathrm{actual}/\mathrm{threshold}, 1) = 1$.  Its
above-sufficiency benchmark provides credit only if
$\mathrm{actual}_k$ genuinely exceeds the threshold, which
re-measurement from a different position does not achieve.
Downstream non-degradation
(Definition~\ref{def:downstream_safe}(b)) is unaffected because
the satisfaction pathway is unchanged.
\end{proof}

\begin{proposition}[S1--Sufficiency Connection]
\label{prop:s1_sufficiency}
Below-sufficiency sacrifices fail Assumption~S1 (free choice,
Section~\ref{sec:sacrifice_model}).  S1 requires that retaining
the sacrificed capability $X_n$ was a \emph{non-degrading}
alternative---one that would not leave the agent below minimum
viable functionality on any need dimension.  For an agent whose
capability level on dimension $d_k$ is below the sufficiency
threshold $s_k$, retaining $X_n$ (not making the trade) leaves
the agent below functionality on $d_k$.  This is not a
non-degrading alternative per S1's definition, so S1 fails for
that event.
\end{proposition}

\noindent\textbf{Consequence for the sacrifice channel.}
Below-sufficiency events are filtered from the revealed-sacrifice
aggregate by S1 before they reach the aggregate lower bound
(Theorem~\ref{thm:aggregate_bound}).  The sacrifice channel
operates only over events where S1 holds.  For above-sufficiency
(want-pursuit) events, the need-insufficiency failure mode of S1
is removed.  S1 may still fail for above-sufficiency events via
other mechanisms (coercion, duress---see
Remark~\ref{rem:s1_above} below).

\begin{remark}[S1 is not automatically satisfied above
sufficiency]
\label{rem:s1_above}
S1 can still fail for above-sufficiency events if the trade was
coerced by other mechanisms (economic duress, political pressure,
information asymmetry) even though the need-sufficiency condition
is met.  The sufficiency threshold resolves one structural failure
mode of S1 (the involuntary-need-satisfaction case); other S1
failure modes (Open Question~\ref{oq:coercion},
Section~\ref{sec:discussion}) remain.
\end{remark}

\begin{proof}[Proof sketch of
Proposition~\ref{prop:s1_sufficiency}]
Let $(i, X, Y, t)$ be a sacrifice event where, absent this trade,
agent~$i$'s capability level on need dimension $d_k$ would fall
below the sufficiency threshold $s_k$.  Agent~$i$ must reach $s_k$
on $d_k$ to maintain functionality.  The alternative of retaining
$X$ (not making the trade) leaves the agent below functionality on
$d_k$.  By S1's non-degrading alternative requirement, retaining
$X$ must not cause the agent to fall below minimum viable
functionality.  But retaining $X$ in this case does exactly that.
Therefore S1 fails for this event, and the event does not enter
the sacrifice aggregate.
\end{proof}

\subsection{How sufficiency thresholds are set}
\label{sec:threshold_setting}

Four mechanisms, not mutually exclusive:

\paragraph{(a) Poset-structural.}
Capabilities with high fan-out (many downstream dependents) have
their sufficiency level set by the minimum required for downstream
capabilities to function.  Computable from poset topology.

\paragraph{(b) Empirical / cross-sectional.}
Observe the population's actual capability levels and identify the
threshold below which agents demonstrably cease functioning on
downstream capabilities.

\paragraph{(c) Revealed-sacrifice-derived.}
Sufficiency is the level at which sacrifice polarity flips.
Below sufficiency, sacrifice is involuntary and increasing; above
sufficiency, sacrifice becomes voluntary and choice-driven.  The
transition point in the sacrifice data identifies the threshold
empirically.  This is the most framework-native option but has a
bootstrap problem: identifying the threshold requires classifying
events, but classifying events requires the threshold.
Resolution: start with structural/normative thresholds
(mechanisms~(a) and~(d)), refine using sacrifice data.

\paragraph{(d) Normative / external.}
A human authority or governance process declares sufficiency
levels.  This is the approach of~\citet{sen1999development}
and~\citet{nussbaum2011creating}.  The framework's formal
apparatus is agnostic to the source; it needs the thresholds as
input.

\medskip\noindent
\textbf{Primary mechanism:} (c) with~(a) as structural
validation.  The revealed-sacrifice data shows where the
transition from involuntary to voluntary sacrifice occurs; the
poset topology confirms the transition corresponds to a
structurally meaningful threshold.
